\documentclass[journal]{IEEEtranTIE}
\usepackage{graphicx}
\usepackage[noadjust]{cite}
\usepackage{amsmath}
\usepackage{url}
\usepackage{amssymb}
\usepackage{color}

\begin{document}
\title{\textcolor{blue}{Semantic Margin-Tightened MPC-SECBF for Context-Dependent Dynamic Obstacle Avoidance}}
    %\title{Social-Aware Extended Euclidean Distance-Based MPC for Context-Dependent Dynamic Obstacle Avoidance}

    \author{Zhang Bo, Lu Xuran
    \thanks{This research was funded in part by the National Natural Science
    Foundation of China under Grant 62173234, in part by the National Major Scientific
    Instruments and Equipments Development Project of National Natural Science Foundation
    of China under Grant 62327808, in part by the Guangdong Major Project of
    Basic and Applied Basic Research under Grant 2023B0303000009 and Grant
    2023B1111050010. \textit{(Corresponding author: Junjie Wen.)}}
    \thanks{Bo Zhang is with the College of Mechatronics and Control Engineering,
    Shenzhen City Joint Laboratory of Autonomous Unmanned Systems and Intelligent
    Manipulation, Shenzhen University, Shenzhen 518060, China and Guangdong
    Laboratory of Artificial Intelligence and Digital Economy (SZ), Shenzhen 518107,
    China. (email: zhangbo@szu.edu.cn).}
    \thanks{Jian Pan and Zihao Huang are with the College of Mechatronics and
    Control Engineering, Shenzhen City Joint Laboratory of Autonomous Unmanned Systems
    and Intelligent Manipulation, Shenzhen 518060, China. (email:
    2210295082@email.szu.edu.cn; 2310295065@email.szu.edu.cn).}
    \thanks{Junjie Wen is with Guangdong Laboratory of Artificial Intelligence
    and Digital Economy (SZ), Shenzhen 518107, China and the Department of
    Mechanical and Automation Engineering, the Chinese University of Hong Kong,
    Shatin N.T., Hong Kong S.A.R. (email: jjwen@mae.cuhk.edu.hk).}
    \thanks{Zongze Wu is with the College of Mechatronics and Control
    Engineering, Shenzhen University, Shenzhen 518060, China and Guangdong Laboratory
    of Artificial Intelligence and Digital Economy (SZ), Shenzhen 518060, China (e-mail:
    zzwu@szu.edu.cn).}
    \thanks{Ben M. Chen is with the Department of Mechanical and Automation Engineering,
    the Chinese University of Hong Kong, Shatin N.T., Hong Kong S.A.R. (e-mail:
    bmchen@cuhk.edu.hk).}
    }
    \maketitle
    \begin{abstract}
\textcolor{blue}{This paper addresses context-dependent dynamic obstacle avoidance for autonomous mobile robots in social dynamic environments, where safety requirements vary with obstacle category and interaction state, including head-on approach, short-term collision urgency, and local crowding. A semantic margin-tightened MPC-ECBF framework is therefore proposed to incorporate category- and interaction-aware safety requirements into velocity-aware dynamic obstacle avoidance through a bounded semantic safety margin. Within this framework, a semantic extended Euclidean safety metric (SEESM) is constructed by preserving the dynamic geometric risk kernel of the extended Euclidean safety metric and tightening it with the bounded semantic safety margin. The semantic margin is then computed from these category and interaction-risk factors. To prevent infeasible or overly conservative boundary tightening, the semantic margin is constrained by category-dependent upper bounds, positive-increment limits, the currently available dynamic geometric safety margin, and an MPC feasibility guard. The resulting SEESM is embedded into a hierarchical planning framework, where it is used for global motion-primitive collision checking and for constructing semantic extended control barrier function constraints in local MPC. A closed-loop practical safety lower bound is derived to quantify the effects of bounded CBF relaxation and semantic-margin updates. Component-level tests, simulation benchmarking, and real-world experiments validate the proposed MPC-SECBF framework in terms of semantic safety clearance, feasibility preservation, and real-time dynamic obstacle avoidance.}   

    \end{abstract}
       % Safe navigation in human-centered dynamic environments requires autonomous mobile robots to account for both motion-induced collision risk and context-dependent safety sensitivity. Conventional Euclidean safety metrics evaluate obstacle avoidance mainly from instantaneous geometric distance, while velocity-aware metrics and MPC-CBF methods still usually treat obstacles as geometric or kinematic entities. This paper proposes a semantic extended Euclidean safety metric (SEESM) for context-dependent dynamic obstacle avoidance. SEESM augments the dynamic geometric safety distance of the extended Euclidean safety metric with a semantic safety margin determined by obstacle category and context-risk factors, including head-on approach, short-term collision urgency, and local crowding. To avoid infeasible or overly conservative boundary tightening, the semantic margin is constrained by category-dependent upper bounds, positive-increment limits, the currently available dynamic geometric safety margin, and an MPC feasibility guard. The proposed metric is embedded into a hierarchical planning framework, where SEESM is used for global motion-primitive collision checking and for constructing semantic extended control barrier function constraints in local MPC. A closed-loop practical safety lower bound is further derived to quantify the effects of CBF relaxation and semantic-margin updates. Component-level tests, simulation benchmarking, and real-world verification are used to evaluate the proposed SEESM-based MPC-SECBF framework for autonomous mobile robot navigation in dynamic environments.
\begin{IEEEkeywords}
Autonomous mobile robots, context-dependent dynamic obstacle avoidance, semantic margin tightening, model predictive control, control barrier function.
\end{IEEEkeywords}
    

    \markboth{IEEE TRANSACTIONS ON INDUSTRIAL ELECTRONICS}


    \section{Introduction}

    \IEEEPARstart{I}{n} \textcolor{blue}{social dynamic environments, avoidance behavior is inherently context-dependent. Humans do not determine safe passing distances solely from geometric proximity; they also consider the type of obstacle and the evolving interaction condition. For example, a fast approaching cyclist, a child in a crowded area, and a static box imply different safety requirements even under similar geometric distances. This motivates context-dependent dynamic obstacle avoidance for autonomous mobile robots, in which safety boundaries should adapt to obstacle category and interaction state.}

\textcolor{blue}{Existing dynamic obstacle avoidance methods have provided effective tools for modeling motion-induced collision risk. Prediction-aware and risk-aware planners reason about future obstacle occupancy, relative motion, and uncertainty to support proactive collision avoidance in dynamic scenes % ~\cite{KEY_RISK_AWARE_PLANNING,KEY_DYNAMIC_SCENE_PREDICTION}.
In parallel, MPC-CBF and dynamic CBF formulations embed safety constraints into predictive trajectory optimization, providing a principled mechanism to balance tracking performance, feasibility, and collision avoidance %~\cite{KEY_DISCRETE_MPC_CBF,KEY_DYNAMIC_CBF_MPC}.
Recent adaptive-margin CBF methods further show that safety margins can be adjusted according to interaction uncertainty or human-related factors%~\cite{KEY_ADAPTIVE_MARGIN_CBF}
. These advances establish a strong foundation for safety-critical navigation in dynamic environments. In social dynamic environments, this foundation can be further extended by allowing the safety boundary to reflect category- and interaction-dependent requirements, rather than being determined only by geometric size and kinematic state.}


\textcolor{blue}{Semantic and social safety-aware navigation has recently become an important direction for extending geometric obstacle avoidance toward human-centered environments. One line of work incorporates semantic risk into spatial safety representations. Semantic labels have been used to inflate occupied regions, reshape distance fields, and tune barrier-function gains, so that semantically sensitive objects induce more conservative safety responses than ordinary obstacles%~\cite{KEY_SEMANTIC_RISK_CBF,KEY_SEMANTIC_SAFETY_FILTER}.
Another line of work constructs semantic safety functions or guidance fields from perception-level semantics. Social-semantic safety filters, Poisson safety functions, Laplace guidance fields, and language-conditioned safety filters translate semantic labels, contextual rules, or natural-language safety specifications into differentiable safety functions or controller constraints%~\cite{KEY_SAFE_SAGE,KEY_POISSON_SAFETY_FUNCTION,KEY_LANGUAGE_SAFETY_FILTER}.
Beyond semantic safety representation, social navigation studies have explored how interaction preferences can be incorporated into planning and control. Learned social zones and proxemic constraints can be encoded as barrier functions or cost terms to enforce human-like clearance, yielding, and personal-space preservation%~\cite{KEY_SOCIAL_ZONE_CBF}.
Topological and homotopy-based planning methods represent passing-side preferences through winding numbers, homotopy classes, or social costs, enabling robots to select socially compliant paths under dynamic pedestrian interactions%~\cite{KEY_TOPOLOGICAL_SOCIAL_MPC,KEY_WINDING_SOCIAL_NAV}.
Recent adaptive-margin, trust-aware, and uncertainty-aware CBF/MPC methods further suggest that safety margins need not be fixed, but can be adjusted according to interaction uncertainty, human-related factors, prediction confidence, or online risk estimates%~\cite{KEY_TRUST_ADAPTIVE_CBF,KEY_CONFORMAL_RISK_CBF}.
These studies demonstrate that semantic categories, social norms, and interaction context can be connected to low-level safety control. Nevertheless, existing formulations usually follow one of three routes: reshaping a spatial safety field, enforcing social zones as hard constraints, or adding social preferences at the planning layer. For velocity-aware dynamic obstacle avoidance, it remains insufficiently explored how category- and interaction-dependent safety requirements can be incorporated into an MPC-ECBF formulation as a bounded tightening of the dynamic safety margin, while preserving the physical interpretation of the EESM dynamic-risk kernel and keeping soft passing preferences separate from hard safety constraints.}

\textcolor{blue}{Motivated by the above observations, this paper proposes a semantic margin-tightened MPC-ECBF framework for context-dependent dynamic obstacle avoidance. The key idea is to preserve the velocity-aware dynamic risk kernel of EESM and to introduce semantic information only through a bounded margin tightening term. Specifically, a semantic safety margin is generated from obstacle category and interaction-risk factors, including head-on approach, short-term collision urgency, and local crowding. The margin is further constrained by category-dependent upper bounds, positive-increment limits, the currently available dynamic geometric safety margin, and an MPC feasibility guard. In this way, semantic risk is incorporated into the near-horizon ECBF constraints without modifying the physical interpretation of the EESM kernel. Side-passing preference is introduced only as a soft cost in the MPC objective, so that it guides socially consistent lateral behavior without changing the hard feasible set.}

\textcolor{blue}{The main contributions of this paper are summarized as follows:
\begin{enumerate}
    \item A semantic margin-tightened MPC-ECBF framework is developed for context-dependent dynamic obstacle avoidance, where category- and interaction-aware safety requirements are incorporated into velocity-aware obstacle avoidance while preserving the EESM dynamic-risk kernel.
    \item A guarded semantic-margin generation and update mechanism is proposed to prevent infeasible or overly conservative boundary tightening by enforcing category-dependent upper bounds, positive-increment limits, available-margin projection, and MPC feasibility checking.
    \item Theoretical properties of the proposed formulation are derived to make the semantic tightening analyzable and compatible with MPC feasibility, including degeneration consistency, margin-shift equivalence, a feasibility sufficient condition, and a closed-loop practical safety lower bound under bounded CBF relaxation and semantic-margin updates. These results characterize the practical safety and feasibility effects of semantic tightening without claiming full recursive feasibility or unconditional safety under arbitrary semantic perception errors.
    \item Component-level tests, simulation benchmarking, and real-world experiments are conducted to evaluate semantic safety clearance, feasibility preservation, and real-time dynamic obstacle avoidance performance.
\end{enumerate}
The remainder of this paper is organized as follows. Section~II presents the proposed semantic margin-tightened MPC-ECBF framework, including the construction of the semantic safety margin, the SEESM-based safety function, the guarded update mechanism, and the MPC-SECBF formulation. Section~III provides the theoretical analysis of the proposed formulation, including feasibility and practical safety properties. Section~IV reports component-level validation, simulation benchmarking, and real-world experiments. Section~V concludes this paper and discusses future work.}
\section{Problem Formulation} \label{sec:problem_formulation}
This work considers context-dependent dynamic obstacle avoidance for
autonomous mobile robots operating in social dynamic environments. At each sampling time $t$, the context-dependent dynamic obstacle avoidance
task is cast as a finite-horizon MPC problem $\mathcal{P}_t$:
\begin{equation}
\begin{aligned}
\min_{\mathbf{x}_t,\mathbf{u}_t,\boldsymbol{\epsilon}_t}\quad
& J_t + \sum_{k=0}^{N_s-1}\sum_{i=1}^{n_t}\rho_{\epsilon}
\epsilon_{i,t+k}^2 \\
\mathrm{s.t.}\quad
& x_{t|t}=x_t, \\
& x_{t+k+1|t}=f_d(x_{t+k|t},u_{t+k|t}), \\
& u_{t+k|t}\in\mathcal{U},\quad x_{t+k|t}\in\mathcal{X}, \\
& 0\le \epsilon_{i,t+k}\le \epsilon_{\max},\\
& h^{\mathrm{ctx}}_i(X_{t+k+1|t};\eta_{i,t})
  \ge (1-\gamma)h^{\mathrm{ctx}}_i(X_{t+k|t};\eta_{i,t}) \\
& \hspace{9.5em}{}-\epsilon_{i,t+k}, 
\end{aligned}
\label{eq:problem_formulation}
\end{equation}
where $k=0,\ldots,N_s-1$ and $i=1,\ldots,n_t$.
Here, $J_t$ penalizes tracking error and control effort, $\mathcal{U}$ and
$\mathcal{X}$ denote the admissible input and state sets, and
$h^{\mathrm{ctx}}_i(\cdot;\eta_{i,t})$ denotes a context-dependent safety
function parameterized by the safety-requirement variable $\eta_{i,t}$.
The variable $\eta_{i,t}$ is used here only to represent the fact that the
required safety boundary may vary with obstacle category and interaction
state; its concrete construction is developed in the following sections.
The relaxation variable $\epsilon_{i,t+k}$ accounts for limited control
authority and abrupt environmental changes, so the resulting safety property
is interpreted as practical safety rather than unconditional hard safety.

The input information includes the robot state, dynamic obstacle positions
and velocities, short-horizon obstacle predictions, semantic categories, and
interaction states. These quantities are assumed to be available from
perception and tracking modules or from experimental annotations. Semantic
recognition, long-term human intent prediction, and social behavior learning
are outside the scope of this paper. \textcolor{red}{Under these assumptions, the research problem is to construct
$h^{\mathrm{ctx}}_i(\cdot;\eta_{i,t})$ and an admissible update rule for
$\eta_{i,t}$, such that category- and interaction-dependent safety
requirements can be incorporated into velocity-aware MPC-ECBF while avoiding
infeasible or overly conservative boundary tightening and maintaining
real-time implementation.}



\section{Main Method}
\subsection{Framework Overview}
Fig.~\ref{fig:overall_framework} illustrates the overall architecture of the proposed
semantic margin-tightened MPC-SECBF framework. At each sampling time, RGB-D
images, point clouds, and IMU data are used to provide obstacle categories,
obstacle states, short-horizon trajectory predictions, and robot odometry.
The obstacle category $c_{i,t}$ is used to determine a category prior margin,
while the predicted obstacle motion is used to evaluate interaction-risk
features, including head-on approach, short-term collision urgency, and local
crowding. These two sources are combined to generate a context-aware candidate
semantic margin.

The candidate margin is then passed through a feasibility-preserving margin
update module. This module enforces category-dependent upper bounds,
positive-increment limits, available-margin constraints, and an MPC feasibility
guard, yielding a bounded and feasible margin. The resulting margin is used to
construct the SEESM safety function by tightening the EESM velocity-aware
safety boundary. In parallel, SLAM and global planning provide the robot state
estimate and a global reference path. The SEESM safety function is embedded as
SECBF constraints in the local MPC optimizer, which tracks the reference path
and generates the control command for the autonomous mobile robot.

    \begin{figure*}[t!]
        \centering
        \includegraphics[width=\textwidth]{系统框图2.png}
        \caption{\textbf{Semantic Margin-Tightened MPC-SECBF Framework}.}
        \label{fig:overall_framework}
    \end{figure*}
The proposed framework instantiates the abstract problem in
Section~\ref{sec:problem_formulation} through three key steps. First, the
context-dependent safety function
$h_i^{\mathrm{ctx}}(\cdot;\eta_{i,t})$ is realized as the SEESM safety
function, which preserves the velocity-aware EESM dynamic-risk kernel while
introducing semantic boundary tightening. Second, the abstract
safety-requirement parameter $\eta_{i,t}$ is specified as the semantic margin
$\beta_{i,t}$, whose value is determined by obstacle category and
interaction-risk factors. Third, the admissible update of $\eta_{i,t}$ is
implemented as a feasibility-preserving semantic-margin update, which bounds
margin variation and prevents infeasible or overly conservative tightening
before the resulting SEESM constraint is embedded into MPC-ECBF.


\subsection{Semantic Extended Euclidean Safety Metric}

\subsection{Feasibility-Preserving Semantic Margin Update}

\subsection{MPC-SECBF Formulation with Soft Passing Preference}
    

    \subsection{Framework Overview}
    \label{sec:Framework_Overview}

    \subsection{Dynamic Perception}
    \label{sec:M_Dynamic_Perception}


    \section{Semantic Extended Euclidean Safety Margin}

    \subsection{Dynamic Geometric Safety Distance}
    In motion planning, the robot and obstacles are commonly abstracted as
    circular regions. Let $R_r$ and $R_o$ denote the robot and obstacle radii,
    respectively. The inflated radius of obstacle $i$ is denoted by
    $R_i=R_r+R_o$, where $i\in\{1,\ldots,n\}$. The conventional Euclidean safety
    margin (CESM) evaluates safety only from the instantaneous geometric distance
    between the robot and obstacles, and defines the safe region as
    \begin{equation}
    SR_{\mathrm{CE}} =
    \{p(t)\mid \|p(t)-p_i^o(t)\|_2>R_d,\ \forall i\},
    \end{equation}
    where $R_d=R_i+R_{\mathrm{safe}}$, and $R_{\mathrm{safe}}$ is the prescribed
    safety threshold. Here, $p(t)=[p_x(t),p_y(t)]^\top$ denotes the robot
    position, $p_i^o(t)$ denotes the position of obstacle $i$ at time $t$, and
    $\|p(t)-p_i^o(t)\|_2$ denotes their Euclidean distance. According to CESM,
    the current state is regarded as unsafe when the relative distance to any
    obstacle is smaller than $R_d$.

    To describe short-horizon interaction risk induced by dynamic obstacles, the
    velocity-extended idea in EESM is introduced. The dynamic geometric term is
    composed of relative position, relative velocity, and forward integration
    time. The relative position and relative velocity between the robot and
    obstacle $i$ are
    \begin{equation}
    l_i=p(t)-p_i^o(t),\qquad v_i=v(t)-v_i^o(t).
    \end{equation}
    The corresponding basic dynamic geometric safe region is
    \begin{equation}
    SR_{\mathrm{EE}} =
    \left\{p(t)\mid
    \|l_i+\tau_i v_i\|_2>R_d,\ \forall i
    \right\}.
    \end{equation}
    where $v(t)=[v_x(t),v_y(t)]^\top$ is the robot velocity, $v_i^o(t)$ is the
    velocity of obstacle $i$, and $\tau_i$ is the forward integration time along
    the relative velocity direction. The vector $l_i+\tau_i v_i$ represents the
    extended geometric displacement under short-term relative motion, and
    $\|l_i+\tau_i v_i\|_2$ quantifies the dynamic geometric distance. The angle
    $\delta_i$ denotes the angle between $l_i$ and $v_i$. Specifically, the
    forward integration time is defined as
    \begin{equation}
    \tau_i=f_r f_v f_T K_e T_i,\qquad K_e\in(0,1),
    \end{equation}
    where
    \begin{equation}
    T_i=\frac{(\|l_i\|_2-R_i)\cos\delta_i}{\|v_i\|_2},
    \end{equation}
    \begin{equation}
    f_r=\mathrm{sign}(-\cos\delta_i),
    \end{equation}
    \begin{equation}
    f_v=\mathrm{sign}\left(
    (l_i^\top v_i)^2+(R_i^2-\|l_i\|_2^2)\|v_i\|_2^2
    \right),
    \end{equation}
    and
    \begin{equation}
    f_T=\mathrm{sign}(T_{\max}-T_i).
    \end{equation}
    Here, $T_i$ estimates the time for the robot to reach the obstacle collision
    boundary; $f_r$ is a relative-direction factor that checks whether the robot
    and obstacle are moving toward each other; $f_v$ is a velocity-obstacle
    factor that checks whether the relative velocity lies in the collision cone;
    $f_T$ is a time-horizon factor that checks whether the potential collision
    lies within the prediction horizon; and $K_e$ relaxes the influence of the
    forward integration time on the safety distance. The indicator function is
    \begin{equation}
    \mathrm{sign}(x)=
    \begin{cases}
    1, & x>0,\\
    0, & x\le 0.
    \end{cases}
    \end{equation}
    Therefore, the dynamic geometric distance and the corresponding safety
    function are
    \begin{equation}
    d_{\mathrm{EE}}(X_t)=\|l_i+\tau_i v_i\|_2,
    \end{equation}
    and
    \begin{equation}
    h_{\mathrm{EE}}(X_t)
    =d_{\mathrm{EE}}(X_t)-R_d
    =\|l_i+\tau_i v_i\|_2-R_d.
    \end{equation}
    When $h_{\mathrm{EE}}(X_t)>0$, the robot remains safe with respect to
    obstacle $i$ in the dynamic geometric sense. When
    $h_{\mathrm{EE}}(X_t)\le0$, the current state is judged to have potential
    collision risk.

    \subsection{Semantic Safety Margin Generation}
    In shared social navigation environments, different obstacle categories have
    different safety sensitivities. For example, children and cyclists usually
    carry higher uncertainty and interaction risk, while non-social objects such
    as boxes require a smaller additional margin. Therefore, a semantic safety
    margin $\beta_{i,t}$ is introduced for each obstacle to tighten the basic
    dynamic geometric boundary.

    The maximum available semantic margin is determined by the semantic category:
    \begin{equation}
    \bar{\beta}_i=\bar{B}(c_{i,t}),
    \end{equation}
    where $c_{i,t}$ is the semantic category of obstacle $i$ at time $t$. The
    category mapping is
    \begin{equation}
    \bar{B}(c_{i,t})=
    \begin{cases}
    0.7\,\mathrm{m}, & c_{i,t}=\mathrm{child},\\
    0.6\,\mathrm{m}, & c_{i,t}=\mathrm{cyclist},\\
    0.5\,\mathrm{m}, & c_{i,t}=\mathrm{vehicle},\\
    0.4\,\mathrm{m}, & c_{i,t}=\mathrm{pedestrian},\\
    0.4\,\mathrm{m}, & c_{i,t}=\mathrm{unknown},\\
    0.1\,\mathrm{m}, & c_{i,t}=\mathrm{box}.
    \end{cases}
    \end{equation}

    Because category alone cannot capture the risk variation within the same
    class, a context-risk modulation mechanism is further introduced. The context
    risk intensity of obstacle $i$ is
    \begin{equation}
    r_i=\mathrm{clip}\left(
    \lambda_h f_{\mathrm{head},i}
    +\lambda_T TTC_{\mathrm{norm},i}
    +\lambda_\rho \rho_{\mathrm{norm},i},0,1
    \right).
    \end{equation}
    where $f_{\mathrm{head},i}$, $TTC_{\mathrm{norm},i}$, and
    $\rho_{\mathrm{norm},i}$ denote the head-on approach degree, short-term
    collision urgency, and local crowding degree, respectively; and $\lambda_h$,
    $\lambda_T$, and $\lambda_\rho$ are the corresponding weights.
    The head-on feature is
    \begin{equation}
    f_{\mathrm{head},i}=\max(0,-\cos\delta_i).
    \end{equation}
    This term is a continuous counterpart of the relative-direction indicator
    $f_r=\mathrm{sign}(-\cos\delta_i)$. The two terms play different roles:
    $f_r$ determines whether to enable the forward integration time $\tau_i$ in
    the dynamic geometric distance, whereas $f_{\mathrm{head},i}$ continuously
    describes the head-on approach degree for semantic-margin modulation.

    To quantify short-term collision urgency, the closing speed is defined as
    \begin{equation}
    v_i^{\mathrm{clo}}=
    \left[
    -\frac{l_i^\top v_i}{\|l_i\|_2+\epsilon_l}
    \right]_+,
    \end{equation}
    where $[a]_+=\max(a,0)$ and $\epsilon_l>0$ avoids division by zero. The
    approximate time to collision is then
    \begin{equation}
    TTC_i=\frac{[\|l_i\|_2-R_i]_+}{v_i^{\mathrm{clo}}+\epsilon_v},
    \end{equation}
    and its normalized risk intensity is
    \begin{equation}
    TTC_{\mathrm{norm},i}=
    \mathrm{clip}\left(1-\frac{TTC_i}{TTC_{\max}},0,1\right).
    \end{equation}
    A smaller $TTC_i$ makes $TTC_{\mathrm{norm},i}$ closer to one, indicating
    higher short-term collision risk, whereas a larger $TTC_i$ drives this term
    toward zero. The $TTC_i$ term is used only for context-risk modulation and
    does not replace $T_i$ or $\tau_i$ in the dynamic geometric safety distance.

    Local density describes the crowding degree around an obstacle. Given a
    neighborhood radius $R_\rho$, the number of neighboring obstacles around
    obstacle $i$ is
    \begin{equation}
    N_i^\rho=
    \sum_{j\neq i}
    \mathbb{I}\left(\|p_j^o(t)-p_i^o(t)\|_2\le R_\rho\right),
    \end{equation}
    \begin{equation}
    \rho_i=\frac{N_i^\rho}{\pi R_\rho^2},\qquad
    \rho_{\mathrm{norm},i}=
    \mathrm{clip}\left(\frac{\rho_i}{\rho_{\max}},0,1\right).
    \end{equation}
    The corresponding local density and normalized local density are denoted by
    $\rho_i$ and $\rho_{\mathrm{norm},i}$, respectively. The weights satisfy
    \begin{equation}
    \lambda_h\ge0,\qquad \lambda_T\ge0,\qquad
    \lambda_\rho\ge0,\qquad
    \lambda_h+\lambda_T+\lambda_\rho=1.
    \end{equation}

    Based on the category upper bound and the context risk, the raw candidate
    semantic margin is
    \begin{equation}
    \tilde{\beta}_{i,t}
    =\bar{\beta}_i
    \left[\mu_{\min}+(1-\mu_{\min})r_i\right],
    \end{equation}
    where $\mu_{\min}\in[0,1]$ is the minimum activation ratio. To avoid abrupt
    tightening of the semantic boundary between adjacent sampling instants, the
    positive rate of change is limited by
    \begin{equation}
    \hat{\beta}_{i,t}
    =
    \min\left\{
    \tilde{\beta}_{i,t},
    \beta_{i,t-1}+\bar{\Delta}_{\beta,i}
    \right\},
    \end{equation}
    which implies
    \begin{equation}
    [\hat{\beta}_{i,t}-\beta_{i,t-1}]_+
    \le \bar{\Delta}_{\beta,i}.
    \end{equation}

    To ensure that the semantic margin does not exceed the available physical
    safety margin, the admissible semantic set is defined as
    \begin{equation}
    \mathcal{B}_i(X_t)=
    \left\{
    \beta\in[0,\bar{\beta}_i]\mid
    h_{\mathrm{EE}}(X_t)-\beta\ge0
    \right\}.
    \end{equation}
    Equivalently,
    \begin{equation}
    \mathcal{B}_i(X_t)
    =
    [0,\bar{\beta}_i]\cap(-\infty,h_{\mathrm{EE}}(X_t)].
    \end{equation}
    When $h_{\mathrm{EE}}(X_t)\ge0$, this set becomes
    \begin{equation}
    \mathcal{B}_i(X_t)
    =
    \left[
    0,\min\{\bar{\beta}_i,h_{\mathrm{EE}}(X_t)\}
    \right].
    \end{equation}
    The rate-limited candidate is projected onto the admissible set:
    \begin{equation}
    \beta_{i,t}^{\mathrm{safe}}
    =
    \Pi_{\mathcal{B}_i(X_t)}(\hat{\beta}_{i,t}).
    \end{equation}
    The semantic margin used by the MPC-SECBF constraints is
    \begin{equation}
    \beta_{i,t}=
    \begin{cases}
    \beta_{i,t}^{\mathrm{safe}},
    & \mathcal{P}_t^{\mathrm{sem}}(\beta_{i,t}^{\mathrm{safe}})
    \ \mathrm{is\ feasible},\\[1mm]
    \Pi_{\mathcal{B}_i(X_t)}(\beta_{i,t-1}),
    & \mathrm{otherwise}.
    \end{cases}
    \end{equation}
    Thus, the final semantic safety margin is not simply the candidate
    $\hat{\beta}_{i,t}$, but is obtained under the category upper bound, context
    risk, positive-increment limit, current safety margin, and MPC feasibility
    guard.



    \subsection{Semantic Extended Euclidean Safety Metric}
    Given $\beta_{i,t}$, the semantic extended Euclidean safety function is
    \begin{equation}
    h_{\mathrm{SEE}}(X_t;\beta_{i,t})
    =
    h_{\mathrm{EE}}(X_t)-\beta_{i,t},
    \end{equation}
    Substituting the definition of $h_{\mathrm{EE}}$ gives
    \begin{equation}
    h_{\mathrm{SEE}}(X_t;\beta_{i,t})
    =
    \|l_i+\tau_i v_i\|_2-R_d-\beta_{i,t}.
    \end{equation}
    The semantic extended Euclidean safe region is
    \begin{equation}
    SR_{\mathrm{SEE}}
    =
    \left\{
    p(t)\mid
    \|l_i+\tau_i v_i\|_2>R_d+\beta_{i,t},\ \forall i
    \right\}.
    \end{equation}
    Thus, SEESM does not directly change the physical radius of the robot or
    obstacle. Instead, it adds a semantic margin $\beta_{i,t}$ to the basic
    dynamic geometric boundary $R_d$. The semantic collision indicator can be
    written as
    \begin{equation}
    F_{\mathrm{SEE}}
    =
    \|l_i+\tau_i v_i\|_2-(R_d+\beta_{i,t}).
    \end{equation}

    \begin{figure}[t!]
        \centering
        \includegraphics[width=\columnwidth]{image_SEESM/SEESM原理图.png}
        \caption{Principle of the proposed SEESM. The semantic margin $\beta_{i,t}$ tightens the dynamic Euclidean safety boundary according to obstacle category and context risk, producing a semantic safety boundary used by both global collision checking and local MPC-SECBF constraints.}
        \label{fig:seesm_principle}
    \end{figure}

    \subsection{SEESM-Based Path Search}
    During primitive expansion, the method samples the control input and duration,
    and then forward-integrates the system to obtain candidate motion primitives.
    The execution cost is jointly determined by the control input $u(t)$ and the
    primitive duration $T$:
    \begin{equation}
    f_e(t)=\int_0^T \|u(t)\|_2^2\,dt+\rho T,
    \end{equation}
    where $\rho$ is a time-penalty weight that suppresses excessively long
    primitives. The heuristic cost $f_h$ is obtained by solving a two-point
    boundary-value problem to estimate the shortest-distance cost from the
    current state to the goal.

    During collision checking, the semantic safety margin is introduced on top
    of the dynamic geometric safety test, resulting in CheckSEESM($\cdot$). For
    candidate primitive node $m$ and predicted obstacle $i$, let $l_{mi}$,
    $v_{mi}$, and $\tau_{mi}$ denote the relative position, relative velocity,
    and forward integration time, respectively. If
    \begin{equation}
    C_{mi}:
    \|l_{mi}+\tau_{mi}v_{mi}\|_2-(R_{d,i}+\beta_{i,t})\le0.
    \end{equation}
    then the primitive node is regarded as unsafe. This criterion indicates that
    path search accounts not only for the relative motion trend of dynamic
    obstacles, but also for the semantic safety margin $\beta_{i,t}$ determined
    by category and context risk. Therefore, unsafe primitives that violate the
    semantic safety constraint can be removed before trajectory optimization.

    \subsection{Semantic Extended Control Barrier Function}
    To incorporate CBF constraints into dynamic obstacle avoidance, SECBF uses
    SEESM to construct the barrier function and introduces slack variables to
    balance safety and feasibility. Let
    $\mathcal{Z}=\{X\in D\mid h(X)\ge0\}$ be the safe set, where
    $h:D\to\mathbb{R}$ is continuously differentiable. For the input-affine
    system $\dot{x}=f(x)+g(x)u$, if there exists an extended class-$\mathcal{K}$
    function $\gamma(\cdot)$, the continuous-time CBF condition can be written as
    \begin{equation}
    \sup_{u\in U}\left[\dot{h}(X,u)\right]\ge -\gamma(h(X)),
    \qquad \gamma\in K_\infty,
    \end{equation}
    with
    \begin{equation}
    \dot{h}(X,u)
    =
    L_fh(X)+L_gh(X)u+
    \frac{\partial h}{\partial p^o}\frac{\partial p^o}{\partial t}.
    \end{equation}
    where $U$ denotes the feasible input set subject to physical and operational
    limits, and $L_fh(X)$ and $L_gh(X)$ are the Lie derivatives of $h(X)$ along
    the drift term $f(x)$ and control input term $g(x)$, respectively. Compared
    with standard DCBF-based CBF design, the additional term
    $\partial h/\partial p^o$ captures the effect of dynamic obstacle-state
    variation on the safety-function derivative.

    To include CBF constraints in a discrete-time controller, the derivative is
    approximated by the one-step difference
    $\dot{h}(X_k)\approx[h(X_{k+1})-h(X_k)]/\Delta t$. The relaxed discrete-time
    CBF recursion is then
    \begin{equation}
    \begin{aligned}
    h(X_{k+1})
    &\ge(1-\gamma)h(X_k)-\epsilon_k,\\
    &\quad 0<\gamma<1,\quad 0\le\epsilon_k\le\epsilon_{\max}.
    \end{aligned}
    \end{equation}
    where $\gamma$ is treated as a scalar for compact discrete-time notation, and
    $\epsilon_k$ is a nonnegative slack variable that relaxes the CBF constraint
    within a bounded range to improve optimization feasibility. Starting from
    $k=0$, recursion gives
    \begin{equation}
    h(X_k)\ge
    (1-\gamma)^kh(X_0)
    -
    \sum_{i=0}^{k-1}(1-\gamma)^{k-1-i}\epsilon_i.
    \end{equation}
    Since $\epsilon_i\le\epsilon_{\max}$,
    \begin{equation}
    \sum_{i=0}^{k-1}(1-\gamma)^{k-1-i}\epsilon_i
    \le
    \epsilon_{\max}\sum_{j=0}^{k-1}(1-\gamma)^j,
    \end{equation}
    and
    \begin{equation}
    \sum_{j=0}^{k-1}(1-\gamma)^j
    =
    \frac{1-(1-\gamma)^k}{\gamma}.
    \end{equation}
    Therefore,
    \begin{equation}
    h(X_k)\ge
    (1-\gamma)^kh(X_0)
    -
    \epsilon_{\max}
    \frac{1-(1-\gamma)^k}{\gamma}.
    \end{equation}
    If $X_0\in\mathcal{Z}$, namely $h(X_0)\ge0$, then
    \begin{equation}
    h(X_k)\ge
    -\epsilon_{\max}
    \frac{1-(1-\gamma)^k}{\gamma}
    \triangleq -\kappa_k,
    \end{equation}
    where $\kappa_k$ is the instantaneous safety-deviation bound at step $k$. As
    $k$ increases, $\kappa_k$ monotonically approaches
    $\kappa_\infty=\epsilon_{\max}/\gamma$. Consequently,
    \begin{equation}
    h(X_k)\ge-\frac{\epsilon_{\max}}{\gamma}.
    \end{equation}
    To keep the practical safety deviation within the prescribed safety
    threshold, the parameters can be selected to satisfy
    $\epsilon_{\max}\le\gamma R_{\mathrm{safe}}$.

    In dynamic obstacle avoidance, the construction of $h(\cdot)$ directly
    determines the effectiveness of the CBF constraint. Considering the predicted
    obstacle trajectory over the MPC horizon, let $k$ denote a discrete prediction
    step. The parameterized state of obstacle $i$ at step $k$ is
    $[p_{x,i}(k),p_{y,i}(k),v_{x,i}(k),v_{y,i}(k),R_i]^\top$. Based on CESM, the
    dynamic geometric safety distance, and SEESM, the safety functions are
    \begin{equation}
    h_{\mathrm{CE}}(X_k)
    =
    \|l_i(k)\|_2-R_i-R_{\mathrm{safe}},
    \quad k=0,1,\ldots,N-1,
    \end{equation}
    \begin{equation}
    \begin{aligned}
    h_{\mathrm{EE}}(X_k)
    &=
    \|l_i(k)+\tau_i v_i(k)\|_2-R_i-R_{\mathrm{safe}},\\
    &\quad k=0,1,\ldots,N-1,
    \end{aligned}
    \end{equation}
    and
    \begin{equation}
    \begin{aligned}
    h_{\mathrm{SEE}}(X_k;\beta_{i,t})
    &=
    \|l_i(k)+\tau_i v_i(k)\|_2
    -R_i-R_{\mathrm{safe}}-\beta_{i,t},\\
    &\quad k=0,1,\ldots,N-1.
    \end{aligned}
    \end{equation}

    \subsection{Model Predictive Control}
    The trajectory tracking problem is formulated as an MPC optimization task
    that minimizes tracking error while preserving safety and feasibility. The
    robot state and input are
    \begin{equation}
    x_t=
    \begin{bmatrix}
    p_x(t)\\p_y(t)\\\theta(t)
    \end{bmatrix},
    \qquad
    u_t=
    \begin{bmatrix}
    v(t)\\\omega(t)
    \end{bmatrix},
    \end{equation}
    where $\theta(t)$ is the heading angle, and $v(t)$ and $\omega(t)$ are the
    linear and angular velocities. The discrete-time kinematic model is
    \begin{equation}
    x_{t+1}=x_t+
    \begin{bmatrix}
    \cos\theta(t)&0\\
    \sin\theta(t)&0\\
    0&1
    \end{bmatrix}
    u_t\Delta t.
    \end{equation}

    Let $U_t=[u_{t|t},u_{t+1|t},\ldots,u_{t+N-1|t}]$ and
    $\Psi_t=\{\epsilon_{i,k}\}$. The semantic MPC-SECBF problem is
    \begin{equation}
    \begin{aligned}
    \mathcal{P}_t^{\mathrm{sem}}:\quad
    \min_{U,\Psi}\quad
    &p(x_{t+N|t})
    +\sum_{k=0}^{N-1}q(x_{t+k|t},u_{t+k|t})\\
    &+\sum_{i=1}^{n}\sum_{k=0}^{N-1}\psi(\epsilon_{i,k})
    +J_t^{\mathrm{side}} ,
    \end{aligned}
    \end{equation}
    where
    \begin{equation}
    J_t^{\mathrm{side}}
    =
    \sum_{i=1}^{n}\sum_{k=0}^{N_s-1}
    w_{i,t}^{\mathrm{side}}\phi(-g_{i,k}^{\mathrm{side}}).
    \end{equation}
    The optimization is subject to
    \begin{equation}
    x_{t+k+1|t}=f(x_{t+k|t},u_{t+k|t}),
    \quad k=0,\ldots,N-1,
    \end{equation}
    \begin{equation}
    x_{t+k|t}\in X_{\mathrm{tra}}^k,\quad
    u_{t+k|t}\in U,\quad k=0,\ldots,N-1,
    \end{equation}
    \begin{equation}
    x_{t|t}=x_t,\qquad x_{t+N|t}\in X_f,
    \end{equation}
    \begin{equation}
    \begin{aligned}
    h_{\mathrm{SEE}}(X_{t+k+1|t};\beta_{i,t})
    &\ge
    (1-\gamma)h_{\mathrm{SEE}}(X_{t+k|t};\beta_{i,t})
    -\epsilon_{i,k},\\
    &\qquad k=0,1,\ldots,K_cN-1,
    \end{aligned}
    \end{equation}
    and
    \begin{equation}
    \begin{aligned}
    h_{\mathrm{CE}}(X_{t+k+1|t})
    &\ge
    (1-\gamma)h_{\mathrm{CE}}(X_{t+k|t})-\epsilon_k,\\
    &\quad k=K_cN,\ldots,N-1.
    \end{aligned}
    \end{equation}
    The cost terms are
    \begin{equation}
    p(x_{t+N|t})=\|x_{t+N|t}-x^d_{t+N|t}\|_P,
    \end{equation}
    \begin{equation}
    \begin{aligned}
    q(x_{t+k|t},u_{t+k|t})
    &=
    \|x_{t+k|t}-x^d_{t+k|t}\|_Q
    +\|u_{t+k|t}\|_R\\
    &\quad
    +\|u_{t+k+1|t}-u_{t+k|t}\|_S,
    \end{aligned}
    \end{equation}
    and
    \begin{equation}
    \psi(\epsilon_{i,k})=\lambda_\epsilon\epsilon_{i,k}^2.
    \end{equation}
    The objective in $\mathcal{P}_t^{\mathrm{sem}}$ consists of terminal tracking
    error, stage tracking and control costs, CBF slack penalties, and the
    side-passing preference cost. The matrices $P$, $Q$, $R$, and $S$ are
    positive definite weights, while $\lambda_\epsilon$ penalizes slack
    variables. The term $J_t^{\mathrm{side}}$ guides more stable lateral passing
    behavior in the presence of semantically sensitive obstacles. Constraints on
    dynamics, state feasibility, control feasibility, initial state, terminal
    set, near-field SECBF, and long-horizon CESM-based DCBF are specified above.
    The stricter SECBF condition is used in the near field because short-term
    safety is critical for real-time collision avoidance, while the longer
    prediction horizon uses the more relaxed CESM-based DCBF condition to reduce
    conservatism under prediction uncertainty.

    \subsection{Closed-Loop Practical Safety}
    The preceding analysis assumes that the semantic safety margin is fixed
    within a single MPC prediction horizon. In closed-loop execution, however,
    $\beta_{i,t}$ is updated at different sampling instants according to
    perception and context risk, thereby changing the threshold of the safety
    function. This section quantifies the practical safety lower bound of the
    semantic-enhanced system.

    \subsubsection{Semantic Safety Set}
    Semantic enhancement cannot create safety by itself; it can only consume the
    dynamic geometric safety margin currently available. For any obstacle $i$,
    define the admissible semantic safety set at state $X_t$ as
    \begin{equation}
    \mathcal{B}_i(X_t)
    =
    \{\beta\in[0,\bar{\beta}_i]:
    h_{\mathrm{SEE}}(X_t;\beta)\ge0\}.
    \end{equation}
    Since
    \begin{equation}
    h_{\mathrm{SEE}}(X_t;\beta)=h_{\mathrm{EE}}(X_t)-\beta,
    \end{equation}
    we have
    \begin{equation}
    h_{\mathrm{SEE}}(X_t;\beta)\ge0
    \Longleftrightarrow
    h_{\mathrm{EE}}(X_t)-\beta\ge0
    \Longleftrightarrow
    \beta\le h_{\mathrm{EE}}(X_t).
    \end{equation}
    Together with $\beta\in[0,\bar{\beta}_i]$, this yields
    \begin{equation}
    \begin{aligned}
    \mathcal{B}_i(X_t)
    &=
    [0,\bar{\beta}_i]\cap(-\infty,h_{\mathrm{EE}}(X_t)]\\
    &=
    [0,\min\{\bar{\beta}_i,h_{\mathrm{EE}}(X_t)\}].
    \end{aligned}
    \end{equation}
    When $h_{\mathrm{EE}}(X_t)\ge0$, this interval at least contains
    $\beta=0$, and thus $\mathcal{B}_i(X_t)\neq\varnothing$. This shows that the
    optional semantic margin is directly bounded by the current dynamic
    geometric safety margin, providing a feasibility guard for semantic updates.

    \subsubsection{Closed-Loop Practical Safety}

    Define the closed-loop safety quantity
    \begin{equation}
    H_{i,t}:=h_{\mathrm{SEE}}(X_t;\beta_{i,t}).
    \end{equation}
    Since $\beta_{i,t}$ is frozen during the MPC solve at time $t$, the first
    executed step satisfies
    \begin{equation}
    h_{\mathrm{SEE}}(X_{t+1};\beta_{i,t})
    \ge
    (1-\gamma)h_{\mathrm{SEE}}(X_t;\beta_{i,t})
    -\epsilon_{i,t}.
    \end{equation}
    Moreover,
    \begin{equation}
    \begin{aligned}
    H_{i,t+1}
    &=h_{\mathrm{SEE}}(X_{t+1};\beta_{i,t+1})\\
    &=h_{\mathrm{SEE}}(X_{t+1};\beta_{i,t})
    -(\beta_{i,t+1}-\beta_{i,t}).
    \end{aligned}
    \end{equation}
    With
    \begin{equation}
    \Delta\beta_i(t):=[\beta_{i,t+1}-\beta_{i,t}]_+,
    \end{equation}
    it follows that
    \begin{equation}
    H_{i,t+1}
    \ge
    (1-\gamma)H_{i,t}-\epsilon_{i,t}-\Delta\beta_i(t).
    \end{equation}
    Expanding the recursion gives
    \begin{equation}
    H_{i,t}
    \ge
    (1-\gamma)^tH_{i,0}
    -
    \sum_{j=0}^{t-1}(1-\gamma)^{t-1-j}
    \left(\epsilon_{i,j}+\Delta\beta_i(j)\right).
    \end{equation}
    If
    \begin{equation}
    0\le\epsilon_{i,t}\le\epsilon_{\max},\qquad
    0\le\Delta\beta_i(t)\le\bar{\Delta}_{\beta,i},
    \end{equation}
    then the closed-loop practical safety lower bound is
    \begin{equation}
    \inf_{t\ge0}H_{i,t}
    \ge
    -\frac{\epsilon_{\max}+\bar{\Delta}_{\beta,i}}{\gamma}.
    \end{equation}
    This result shows that cross-sampling semantic-margin updates consume
    additional safety margin, and their influence is bounded by the maximum
    positive increment $\bar{\Delta}_{\beta,i}$. A smaller
    $\bar{\Delta}_{\beta,i}$ produces smoother semantic-boundary changes and a
    tighter closed-loop safety lower bound.
    When $\beta_{i,t}\equiv0$ and $w_{i,t}^{\mathrm{side}}\equiv0$, the proposed
    semantic framework degenerates to the baseline hierarchical MPC-ECBF
    structure without semantic terms, and the bound becomes
    \begin{equation}
    \inf_{t\ge0}h_{\mathrm{EE}}(X_t)
    \ge
    -\frac{\epsilon_{\max}}{\gamma}.
    \end{equation}
    \section{Experimental Validation}
    \label{sec:Experiments}

    \subsection{Verification Settings and Evaluation Metrics}
    \label{sec:Verification_Settings}
    The experimental validation is designed to answer three questions. First, it
    verifies whether the proposed semantic margin produces reasonable
    category- and context-dependent safety boundaries. Second, it evaluates
    whether the guarded semantic-margin update is necessary for maintaining MPC
    feasibility and smooth safety-boundary variation. Third, it examines whether
    the complete SEESM-based MPC-SECBF framework improves safety while preserving
    planning efficiency and real-time feasibility in dynamic social navigation
    scenarios.

    All simulation experiments will be conducted in a C++/ROS or ROS/Gazebo
    environment using the same robot model, start and goal positions, obstacle
    trajectories, prediction horizon, MPC parameters, and computing platform for
    all compared methods. Unless otherwise stated, each method will be evaluated
    over 100 repeated trials in each simulation scenario. The real-world
    experiments will be conducted on a Scout 2.0 or equivalent autonomous mobile
    robot platform equipped with onboard LiDAR, IMU, and computing resources.
    Each real-world method configuration will be repeated at least 10 times
    under the same task setting.

    The following metrics are used throughout the experiments. The success rate
    $R_{\mathrm{suc}}$ is the percentage of trials completed without collision
    or deadlock. The minimum obstacle distance $D_{\min}$ measures the smallest
    robot-obstacle separation during navigation. The semantic clearance
    $D_{\mathrm{sem}}$ is defined as the minimum value of
    $h_{\mathrm{SEE}}(X_t;\beta_{i,t})$ over all obstacles and sampling instants,
    and is used to evaluate whether the semantic safety boundary is satisfied.
    The path length $L_{\mathrm{path}}$ measures navigation efficiency. The
    control effort $E_{\mathrm{cost}}$ is computed by integrating the squared
    control input along the executed trajectory. The velocity variation
    $\Phi_{v|\omega}$ evaluates trajectory smoothness from linear and angular
    velocity changes. The average computation time $T_{\mathrm{comp}}$ measures
    planning and optimization efficiency. The MPC feasibility rate, average CBF
    slack, and semantic-margin increment $\Delta\beta$ are further reported to
    quantify the influence of semantic constraints on optimization feasibility
    and boundary smoothness.

    \subsection{Component-Level Validation}
    \label{sec:Component-level-validation}
    Component-level validation isolates the main modules of SEESM before testing
    the full navigation system. These experiments focus on semantic-margin
    generation, guarded margin updates, and closed-loop practical safety.

    \subsubsection{Semantic Margin Generation Test}
    This test verifies whether the semantic safety margin follows the intended
    category and context ordering. Six obstacle categories are considered:
    child, cyclist, vehicle, pedestrian, unknown object, and box. For each
    category, four interaction contexts are constructed: head-on approach,
    crossing motion, moving-away motion, and crowded local interaction. The
    generated semantic margin $\beta_{i,t}$ is recorded together with the
    category upper bound $\bar{\beta}_i$, head-on factor
    $f_{\mathrm{head},i}$, normalized time-to-collision
    $TTC_{\mathrm{norm},i}$, and normalized local density
    $\rho_{\mathrm{norm},i}$.

    The expected evidence is that semantically sensitive objects, such as
    children and cyclists, receive larger margins than non-social objects under
    comparable geometric conditions, and that the margin increases under
    head-on, short-TTC, and crowded contexts. The result should be reported as a
    category-context matrix of $\beta_{i,t}$ and as time-series curves showing
    how the margin changes when the interaction context changes.

    \subsubsection{Guard Mechanism Ablation}
    This ablation evaluates whether the guarded update mechanism is necessary
    for feasibility and smooth safety-boundary evolution. The following variants
    are compared: full SEESM, category-only margin, context-only margin, SEESM
    without positive-increment limiting, SEESM without admissible-set projection,
    and SEESM without the MPC feasibility guard. All variants use the same
    perception input, obstacle trajectories, MPC horizon, and cost weights.

    The reported quantities include the maximum positive margin increment
    $\max_t\Delta\beta_i(t)$, the MPC infeasibility rate, the average CBF slack,
    and the minimum closed-loop semantic safety quantity $H_{i,t}$. A useful
    ablation outcome should show whether removing rate limits causes abrupt
    semantic-boundary changes, whether removing admissible-set projection allows
    $\beta_{i,t}$ to exceed the available dynamic geometric margin, and whether
    removing the feasibility guard increases solver infeasibility.

    \subsubsection{Closed-Loop Practical Safety Check}
    This test validates the practical safety bound derived in the closed-loop
    analysis. During closed-loop execution, the quantity
    $H_{i,t}=h_{\mathrm{SEE}}(X_t;\beta_{i,t})$ is recorded for each obstacle and
    compared with the theoretical lower bound
    $-(\epsilon_{\max}+\bar{\Delta}_{\beta,i})/\gamma$. The positive semantic
    increment limit $\bar{\Delta}_{\beta,i}$ is varied over several values while
    keeping the same task and obstacle trajectories.

    The expected presentation is a set of time-series plots for $H_{i,t}$,
    $\beta_{i,t}$, $\Delta\beta_i(t)$, and CBF slack. The experiment should
    report whether the measured minimum of $H_{i,t}$ remains above the derived
    bound and whether smaller $\bar{\Delta}_{\beta,i}$ produces smoother
    semantic-boundary variation and a tighter empirical safety margin.

    \subsection{Simulation Scenario Benchmarking}
    \label{E_simulation-section}

    \subsubsection{Experimental Setup}
    Simulation benchmarking evaluates the complete SEESM-based MPC-SECBF
    framework against representative baselines. The compared methods are:
    CESM-MPC-CBF, which uses only conventional Euclidean distance constraints;
    MPC-DCBF, which considers predicted dynamic obstacle states; EESM-MPC-ECBF,
    which uses the velocity-aware EESM without semantic margins; SEESM without
    the full guard mechanism; and the proposed full SEESM-MPC-SECBF.

    Four simulation scenarios are used. Scenario 1 contains a single head-on
    pedestrian or cyclist and evaluates whether semantically sensitive obstacles
    trigger earlier and smoother avoidance. Scenario 2 contains crossing
    obstacles with mixed categories, including pedestrian, cyclist, and box, and
    tests whether different semantic objects induce different safety boundaries
    under comparable geometric risk. Scenario 3 is a crowded dynamic scene with
    three to six moving obstacles, including child, unknown object, and vehicle,
    and evaluates the effect of local crowding on safety and feasibility.
    Scenario 4 is a stress test with higher obstacle speed, higher obstacle
    density, or shorter time-to-collision, and is used to compare the feasibility
    and safety robustness of full SEESM against unguarded semantic constraints.

    For each scenario, the main table should report
    $R_{\mathrm{suc}}\uparrow$, $D_{\min}\uparrow$ (m),
    $D_{\mathrm{sem}}\uparrow$, $L_{\mathrm{path}}\downarrow$ (m),
    $\Phi_{v|\omega}\downarrow$, $T_{\mathrm{comp}}\downarrow$ (ms), and MPC
    feasibility rate $\uparrow$ (\%). Representative trajectory plots should
    compare the geometric paths of all methods. Time-series plots should show
    $\beta_{i,t}$, $H_{i,t}$, MPC slack, and computation time for the most
    challenging scenario.

    \subsubsection{Results and Discussion}
    The simulation results should be discussed according to the evidence needed
    for each claim. If full SEESM obtains a larger $D_{\mathrm{sem}}$ and
    $D_{\min}$ than EESM-MPC-ECBF under mixed-category scenarios, the result
    supports the value of semantic boundary modulation. If full SEESM maintains a
    higher MPC feasibility rate and lower slack than the unguarded SEESM variant,
    the result supports the necessity of the guarded semantic-margin update. If
    full SEESM maintains $T_{\mathrm{comp}}$ within the control-cycle deadline,
    the result supports real-time applicability.

    Failure cases should also be reported. CESM-MPC-CBF is expected to fail or
    become overly aggressive when relative velocity and short-TTC interactions
    dominate the risk. Unguarded SEESM is expected to become overly conservative
    or infeasible when semantic margins increase abruptly. These cases should be
    illustrated with trajectory snapshots, margin time-series, and infeasibility
    records rather than only with aggregate metrics.

    \subsection{Real-World Verification}
    \label{E_real_world-section}
    \subsubsection{Experimental Setup}
    Real-world verification tests whether the proposed planner can run on a
    physical AMR and produce executable semantic-aware avoidance behavior. The
    platform is a Scout 2.0 or equivalent differential-drive AMR equipped with
    onboard LiDAR, IMU, and a computing unit. The robot performs approximately
    20 m navigation tasks in outdoor or indoor dynamic environments.

    Three real-world scenarios are recommended: a single pedestrian crossing the
    robot path, a mixed pedestrian-cyclist passing scenario, and a pedestrian-box
    mixed scenario. The semantic category can be provided by an online semantic
    detector if available. If semantic recognition is not part of the contribution,
    the categories can be provided by manual annotation or an experimental script;
    in that case, the paper should state explicitly that the experiment evaluates
    the SEESM planner rather than semantic classification accuracy.

    At least two methods should be compared in the real-world tests:
    EESM-MPC-ECBF and the proposed full SEESM-MPC-SECBF. If safety conditions
    allow, an unguarded SEESM variant can be tested at low speed with a safety
    operator, but it should not be used in risky close-interaction trials. The
    reported metrics are $D_{\min}$, $D_{\mathrm{sem}}$, $L_{\mathrm{path}}$,
    $E_{\mathrm{cost}}$, $\Phi_{v|\omega}$, success rate, and average control
    frequency.

    \subsubsection{Results and Discussion}
    The real-world results should focus on whether SEESM changes the robot's
    behavior in a semantically meaningful and executable way. In the pedestrian
    and cyclist scenarios, the proposed method should be evaluated by whether it
    increases semantic clearance and avoids abrupt control changes. In the
    pedestrian-box scenario, the result should show whether the robot maintains a
    larger margin for the social obstacle while avoiding unnecessary detours
    around the non-social object.

    The final presentation should include trajectory overlays, selected
    snapshots of the robot-obstacle interaction, and a quantitative table. The
    discussion should remain bounded: real-world verification demonstrates
    platform feasibility and semantic-aware behavior under the tested scenarios,
    while broader generalization to all social navigation environments requires
    additional evaluation with more diverse semantic perception inputs.


    \section{Conclusion}
    \label{sec:Conclusion}


    \section*{Acknowledgment}

    \bibliographystyle{Bibliography/IEEEtranTIE}
    \bibliography{Bibliography/IEEEabrv,Bibliography/BIB_xx-TIE-xxxx}

\end{document}
